\section{Compiling a Fourier transform on its own coordinates}
\label{loc:section}
\label{sec:local}

We next express a Fourier matrix using shallow layers on exactly its
given coordinates.  Preserving the width is essential: later, the
coordinate sets will be tensor factors, so padding them separately would
multiply the total dimension.  We first factor the Fourier matrix into
diagonal and triangular Toeplitz matrices.  We then implement those Toeplitz factors on their
own coordinates, using the exact-width mechanism of
\cite[Section~3]{exactfourier} with an effective choice of workspace.
Finally, we compile every two-coordinate operation to the particular
matrix $C$ used above.

A \emph{shear} adds a scalar multiple of one coordinate to a distinct
coordinate and leaves all other coordinates unchanged.  A \emph{shear
layer} consists of shears on disjoint pairs.  A \emph{monomial layer} is
an invertible monomial matrix.  We also use \emph{mixed rounds}: on
disjoint coordinate regions, such a round performs monomial layers or
shear layers.  It splits into a monomial layer followed by a shear
layer.  This convention lets independent recursive calls run in parallel.

\begin{proposition}[Uniform local compiler]\label{loc:compiler}
Let $r\ge1$ and let $u$ be any primitive $r$th root of unity.  The matrix
$F_r(u)=(u^{ij})_{0\le i,j<r}$ has a word on exactly $r$ coordinates
with $O(1+\log^4r)$ slots, each of which is either
\begin{enumerate}
\item a monomial layer, or
\item disjoint copies of $C$ on ordered pairs, with identity on every
remaining singleton.
\end{enumerate}
There is one fixed repeating pattern of slot types into which all these
words can be padded.  For any finite collection of widths at most $R$, the words
can therefore be synchronized in $O(1+\log^4R)$ slots.

Their schedules and scalar coefficients can be prepared in time
$r^{O(1)}$, with an absolute exponent, counting both exact scalar
arithmetic and integer bookkeeping.  Besides $u$ and the fixed constant
$i$, this preparation needs only power-of-two roots of orders at most
$8r$.  Coefficients are specified by shared arithmetic programs over
these roots and rational numbers.  Every division is by a proved
nonzero scalar; every structural choice uses integers alone.  Neither
a test for zero of a complex number nor conjugation of input data is
required.
\end{proposition}

\subsection{Reducing Fourier matrices to triangular Toeplitz matrices}

The Fourier matrix evaluates a polynomial at the powers of $u$.
Writing that polynomial in the Newton basis makes evaluation triangular.
Symmetry then removes the need to compile the change of basis separately.

\begin{lemma}[Newton reduction]\label{loc:newton-reduction}
For every primitive $r$th root $u$, with $r\ge1$, there is an invertible
lower triangular matrix $N$ such that
\[
 F_r(u)=N D_N^{-1}N^{\mathsf T},
 \qquad D_N=\diag(N_{00},\ldots,N_{r-1,r-1}).
\]
The matrix $N$ is a product of two invertible diagonal matrices and
one invertible lower triangular Toeplitz matrix of width $r$.
All these factors can be prepared in polynomial time as rational
expressions in $u$, with every denominator nonzero.
\end{lemma}

\begin{proof}
For $r=1$ take $N=D_N=T=(1)$.  Otherwise let the $k$th column of $P$
be the coefficient vector of the monic Newton polynomial
$\prod_{j=0}^{k-1}(Z-u^j)$, for $0\le k<r$.  Then $P$ is unit upper
triangular and $N=F_r(u)P$ is lower triangular.  For $i\ge k$,
\begin{equation}\label{loc:newton}
 N_{ik}=\prod_{j=0}^{k-1}(u^i-u^j)
       =(-1)^ku^{k(k-1)/2}\frac{H_i}{H_{i-k}},
 \qquad H_j=\prod_{s=1}^j(1-u^s),\quad H_0=1.
\end{equation}
The required $H_j$, $0\le j<r$, are nonzero because $u$ is primitive.
Consequently
\[
 N=\diag(H_i)\,T\,
       \diag\bigl((-1)^ku^{k(k-1)/2}\bigr),
 \qquad
 T_{ik}=\begin{cases}H_{i-k}^{-1},&i\ge k,\\0,&i<k.\end{cases}
\]
Here $T$ is invertible lower triangular Toeplitz, with diagonal one
and coefficient sequence $h_j=H_j^{-1}$.

Let $D_N=\diag(N_{00},\ldots,N_{r-1,r-1})$ and $L=ND_N^{-1}$.
The identity $F_r(u)=L D_NP^{-1}$ is a unit-lower/diagonal/unit-upper
factorization.  Such a factorization is unique: comparing two makes
a unit lower triangular matrix equal an upper triangular matrix,
forcing both to be identity, and then identifies their diagonal
factors.  Since $F_r(u)=F_r(u)^{\mathsf T}$, transposition gives another
such factorization.  Uniqueness yields $P^{-1}=L^{\mathsf T}$, hence
\begin{equation}\label{loc:fourier-factorization}
 F_r(u)=N D_N^{-1}N^{\mathsf T}.
\end{equation}
Computing the powers of $u$, the $H_j$, and all diagonal entries has
polynomial cost.  Every division above is by a product of nonzero
$H_j$ and powers of $u$.  In particular, the argument applies to
composite $r$ and to every primitive root of order $r$.
\end{proof}

Thus it suffices to compile an invertible triangular Toeplitz matrix
on its own coordinates.  Its transpose will require the same number
of layers: transposition reverses the factor order and preserves
monomial and shear layers.  The next two subsections establish this
Toeplitz compiler, including the workspace needed by its convolutions.

\subsection{Convolution and arbitrary borrowed coordinates}

We use linear arithmetic directed acyclic graphs (DAGs) in which each
gate forms a linear combination of at most two sources or earlier gate
values.  Each output is one designated source or gate read, or zero;
any output sum or scaling is included among the gates.  Fan-out counts
these designated output uses as well as uses by later gates.

The elementary convolution circuits will supply the local operations.
A convolution of variable data with a fixed vector, whose operand lengths have sum $v$, has a linear
arithmetic DAG with
\begin{equation}\label{loc:convolution-bounds}
 O\bigl(v\log(2+v)\bigr)\text{ gates},\qquad
 O\bigl(\log(2+v)\bigr)\text{ depth},
\end{equation}
and absolute bounded fan-out, counting designated output uses.  Gates
have at most two predecessors.  Reversals and truncations preserve
these bounds.

Indeed, pad to the least power of two $L$ at least the full product
length.  At $L$th roots of unity, convolution is pointwise
multiplication.  The transform of the fixed vector supplies scalar
coefficients.  The binary Fourier recursion separates even and odd
coefficients and combines transformed entries $a,b$ into
$a+\omega b,a-\omega b$.  Each entry is used a bounded number of times
at the next level.  The inverse uses inverse roots and multiplication
by $L^{-1}$, as follows from the geometric sum of the roots.  This
proves \eqref{loc:convolution-bounds}, including all multiplications by
fixed scalars.  Literal zero inputs need no variable registers.
Selections have no repeated output indices; a fixed number of these
circuits, their compositions, and their sums still have absolute
bounded fan-out.  Their topology can be printed before any scalar
coefficients are evaluated, retaining even gates with a coefficient
that might happen to vanish.

The DAG may have extra wires. The following identity implements its
effect using coordinates that already carry other data. Its first replay
has the compute--copy--reverse pattern of Bennett's reversible
simulation~\cite{bennett1973}. Restoring storage with arbitrary initial
contents is also the principle of clean transparent computation in
catalytic space~\cite[Section~3]{buhrmanEtAl2014}. The source-suppressed second
replay below proves the exact cancellation and layer bounds needed here.

\begin{lemma}[Replay with arbitrary borrowed coordinates]\label{loc:replay}
Suppose $M:\C^e\to\C^a$ has a linear DAG with $g$ gates, depth $H$,
fan-in at most two, and fan-out at most $\delta$, including output
uses.  Each output is a source or gate read, or zero.  On disjoint
coordinate sets $x,y,z$ of sizes $e,a,g$, respectively,
a word of $O((\delta+1)(H+1))$ shear layers implements
\[
 (x,y,z)\longmapsto(x,y+Mx,z).
\]
It uses at most $8g+2a$ shears.  The initial contents of $z$ are arbitrary,
including values correlated with $x$, $y$, or other data.
\end{lemma}

\begin{proof}
Assign one coordinate $z_v$ to each gate, in topological order.  Instead
of assigning the gate value, add its prescribed linear combination of
predecessors into $z_v$.  This takes at most two shears.  The resulting
forward sweep has the form
\[
 z\longmapsto Uz+Tx,
\]
where $U$ is unit lower triangular.  Write the output reads as $Fz+Dx$,
allowing outputs that read a source directly; each nonzero output is
one designated read.  Starting with zero gate coordinates computes the
original DAG, so $FT+D=M$.

Run the forward sweep, add its output reads into $y$, and undo the
sweep.  This restores $z$ and adds $Mx+FUz$ to $y$.  Next run the same
sweep with all direct-source edges omitted, omit $Dx$ at the readout,
subtract the resulting output $FUz$, and undo the sweep again.  The net
effect is the stated map.  The target coordinates are never
predecessors in either sweep.  The calculation is an identity in
independent variables $x,y,z$, so it remains valid for correlated
initial contents.

To arrange a sweep into layers, assign gates their longest-path depth
levels.  Within one level, destinations are new gate coordinates and
predecessors are earlier gate coordinates or sources.  The undirected
multigraph of shear edges has maximum degree
$\Delta\le\max\{2,\delta\}$.  Greedy edge coloring uses at most
$2\Delta-1$ colors, since each edge meets at most $2\Delta-2$ others.
The operations in one color have disjoint pairs.  Reordering by color
is valid because no destination is a predecessor at the same level.
The output additions have the same property: their sources have
bounded output fan-out and each target has one read.  Reverse sweeps
use the reversed layers and negated coefficients.  Four sweeps and
two output rounds give the depth bound and the count $8g+2a$.
\end{proof}

\subsection{Toeplitz layers and a measured workspace rule}

The matrix representation in the next proof is the displacement-rank
construction of Kailath, Kung, and Morf~\cite[Lemmas~1--2]{kailathKungMorf1979}:
low-rank shift differences reconstruct a matrix through Toeplitz
convolution factors. We give the rectangular-chunk formulas and then
fit their computation into the available coordinates.

\begin{lemma}[Effective exact-width Toeplitz layers]\label{loc:toeplitz}
An invertible lower triangular Toeplitz matrix of width $r$ has a word
on exactly $r$ coordinates in $O(1+\log^4r)$ mixed rounds.  A fixed
algorithm prepares the word and its coefficients in $r^{O(1)}$ time
from the Toeplitz entries and power-of-two roots of orders at most
$8r$.  The topology depends only on $r$; all divisions by Toeplitz
parameters are justified by the nonzero diagonal entry.
\end{lemma}

\begin{proof}
We split the transform into two smaller diagonal blocks and a cross
update.  For each rectangular part of that update, we will construct a
convolution DAG and fit all its gate coordinates inside the parent
block.  We then count the resulting layers and the preparation work.

Write $T_v=(h_{i-j})_{0\le j\le i<v}$, with zero entries above the
diagonal and $h_0\ne0$.  The same initial sequence $h$ specifies all
smaller blocks used below.  Its reciprocal-series coefficients are
computed by
\begin{equation}\label{loc:reciprocal}
 g_0=h_0^{-1},\qquad
 g_k=-h_0^{-1}\sum_{j=1}^k h_jg_{k-j}\quad(1\le k<r).
\end{equation}
They give the inverse of every lower triangular Toeplitz block.

For $v>1$, put $p=\lfloor v/2\rfloor$ and write
\[
 T_v=\begin{pmatrix}T_p&0\\B&T_{v-p}\end{pmatrix}.
\]
The factorization
\[
 T_v=\begin{pmatrix}I&0\\E&I\end{pmatrix}
      \diag(T_p,T_{v-p}),\qquad E=BT_p^{-1},
\]
first performs the two diagonal transforms in parallel, then adds to
the second block from the transformed first block.  In the original
indices of this cross block,
\begin{equation}\label{loc:cross}
 E_{ij}=\sum_{\nu=j}^{p-1}h_{i-\nu}g_{\nu-j},
 \qquad p\le i<v,\quad 0\le j<p.
\end{equation}

Partition both intervals into consecutive chunks.  For a target chunk
of size $a$ and source chunk of size $e$, let $M$ be the corresponding
$a\times e$ submatrix of $E$.  Let $J_a,J_e$ denote the lower shifts.
Away from the first local row and column, \eqref{loc:cross} telescopes
to
\begin{equation}\label{loc:displacement-entry}
 E_{ij}-E_{i-1,j-1}=-h_{i-p}g_{p-j}.
\end{equation}
Thus, writing $\mathcal D=M-J_aMJ_e^{\mathsf T}$, we can form
\begin{equation}\label{loc:displacement}
 \mathcal D=\sum_{\nu=1}^3 v^{(\nu)}(w^{(\nu)})^{\mathsf T}
\end{equation}
without any rank test.  Specifically, extend the outer product on the
right of \eqref{loc:displacement-entry} to the entire chunk and subtract
it from $\mathcal D$, obtaining $R$.  If $e_0$ is the first unit vector,
the two remaining outer products are
$e_0R_{0,*}$ and $(R_{*,0}-R_{00}e_0)e_0^{\mathsf T}$.
They correct its first row and first column.  Keep all three terms,
even when some coefficients vanish.

Iterating $M=\mathcal D+J_aMJ_e^{\mathsf T}$ gives, by nilpotence,
the finite reconstruction
\begin{equation}\label{loc:reconstruction}
 M=\sum_{l\ge0}J_a^l\mathcal D(J_e^{\mathsf T})^l.
\end{equation}
For one outer product $vw^{\mathsf T}$, its contribution has entries
$\sum_{l\ge0}v_{i-l}w_{j-l}$, with zero extension.  Apply it to $x$ by
\begin{equation}\label{loc:two-convolutions}
 b_l=\sum_{j=l}^{e-1}w_{j-l}x_j\quad(0\le l<e),\qquad
 c_i=\sum_{l=0}^{e-1}v_{i-l}b_l\quad(0\le i<a).
\end{equation}
If $x^{\mathrm{rev}}_j=x_{e-1-j}$, the first map selects
$b_l=(w*x^{\mathrm{rev}})_{e-1-l}$ from an ordinary convolution.
The second map is a truncated convolution.  Using three
branches and summing their outputs gives a specific DAG with absolute
constants $K_0,K_1,\delta_0$ such that
\begin{equation}\label{loc:chunk-bounds}
 g\le K_0(a+e)\log_2(2+a+e),\qquad
 H\le K_1\log_2(2+a+e),\qquad \delta\le\delta_0.
\end{equation}
These include output uses: every source is used in only three
branches, the branches retain bounded fan-out, and their sum uses a
bounded number of gates per output.

To fit this DAG inside the parent block, reserve the $a+e$ target and
source coordinates and one borrowed coordinate for each of its $g$
gates.  All $g$ coordinates stay assigned to their gates throughout
the replay.  Here is an effective choice of chunk size that does not
require knowing $K_0$.  For every integer $b\in\{1,\ldots,v\}$, partition
each interval into chunks of size $b$, except for its possible last chunk.
Print the DAG skeleton for every resulting pair and count its gates.
Declare $b$ admissible precisely when, for every pair,
\begin{equation}\label{loc:fit}
 g+a+e\le v.
\end{equation}
Choose the largest admissible $b$.  This is an integer test on printed
graphs, independent of the coefficients, made before any data are
transformed.  If no $b$ is admissible, replace the entire recursive
decomposition of $T_v$ by a direct word: process rows in decreasing
order, multiply the current coordinate by $h_0$, and add $h_{i-j}$
times each lower-indexed coordinate $j<i$.  The lower coordinates
still have their original values, so this computes exactly $T_v$.
Width one is a scaling.

To prove the depth bound, fix for analysis a sufficiently large
absolute $D$, for example $D\ge16K_0+16$.  For all sufficiently large
$v$,
\[
 b_0=\left\lfloor\frac{v}{D\log_2v}\right\rfloor
 \ge\frac{v}{2D\log_2v}\ge1
\]
passes \eqref{loc:fit}.  Indeed $a+e\le2v/(D\log_2v)\le v/8$, while
\eqref{loc:chunk-bounds} gives $g\le4K_0v/D\le v/4$.
Consequently the direct fallback occurs only for bounded $v$, and the
chosen size $b\ge b_0$ leaves $O(\log v)$ chunks in each interval.

For a selected pair, borrow $g$ coordinates outside its source and
target chunks but inside these same $v$ coordinates; \eqref{loc:fit}
guarantees their availability.  Lemma~\ref{loc:replay} implements its
update in $O(1+\log v)$ shear layers.  Process the $O(\log^2v)$ pairs
serially.  Every borrowed coordinate is restored before the next pair,
so it may carry a source value, an already updated target value, or any
other current data.  The complete cross update takes
$O(1+\log^3v)$ layers.  Child calls use disjoint blocks, hence the
mixed-round depth satisfies
\[
 d(v)\le\max\{d(\lfloor v/2\rfloor),d(\lceil v/2\rceil)\}
          +O(1+\log^3v).
\]
Summing along a balanced recursion path proves
$d(r)=O(1+\log^4r)$.

We finally account for preparation.  Formula~\eqref{loc:reciprocal},
the entries \eqref{loc:cross}, and the three explicit outer products
use polynomially many scalar operations.  Each node of the recursion
tests at most $v$ sizes and at most $v^2$ chunk pairs per size; every
printed graph has $O(v\log(2+v))$ gates.  There are fewer than $2r$
nodes.  Even this deliberately redundant search has polynomial
integer cost.  Gate-level assignments, bounded-degree edge colorings,
and the choice of the first available borrowed coordinates can all be
performed by finite scans of these lists.  Only the chosen graphs
need their scalar coefficients evaluated.  Fixed-vector transforms
and all other coefficient computations have polynomial cost as well.
Store computed scalars and references to them as an arithmetic DAG;
there is no expansion into expression trees.

The first convolution in \eqref{loc:two-convolutions} has operand
lengths $e,e$, and the second has lengths $a,e$.  Their padded Fourier
lengths are powers of two smaller than $4r$.  Thus a single primitive
root of power-of-two order $L_r=2^{\lceil\log_2(4r)\rceil}<8r$
supplies every needed root by taking powers.  These powers and their
inverses are also prepared with polynomially many operations.  Integer
indices in the entire compiler have $O(\log(2+r))$ bits.  This proves
the claimed preparation bound, with no need to determine the finite
threshold in the depth argument.
\end{proof}

\subsection{Fourier layers and a fixed word in the kernel}

\begin{lemma}[Fourier layers at every primitive root]\label{loc:fourier}
For every primitive $r$th root $u$, the matrix $F_r(u)$ has a
polynomially preparable exact-width word in
$O(1+\log^4r)$ monomial and shear layers.  Its coefficients are rational
expressions in $u$ and the power-of-two roots from
Lemma~\ref{loc:toeplitz}, with all denominators nonzero.
\end{lemma}

\begin{proof}
Apply Lemma~\ref{loc:toeplitz} to the Toeplitz factor of $N$ in
Lemma~\ref{loc:newton-reduction}.  The two diagonal factors of $N$
are monomial layers.  Equation~\eqref{loc:fourier-factorization}
then gives the Fourier word by composing the words for $N^{\mathsf T}$,
$D_N^{-1}$, and $N$.  Transposition preserves the layer types and
reverses the factor order.  Split mixed rounds into monomial and shear
layers.  There are $O(1+\log^4r)$ layers in total, and all preparation
bounds and nonzero-denominator guarantees follow from the two lemmas.
\end{proof}

We complete the proof of Proposition~\ref{loc:compiler}.  Recall
$C=\left(\begin{smallmatrix}a&b\\b&a\end{smallmatrix}\right)$,
$a=(1+i)/2$, $b=(1-i)/2$.  Set
\[
 H'=\begin{pmatrix}1&1\\1&-1\end{pmatrix},\qquad
 S=\diag(1,i).
\]
Direct multiplication gives $H'=a^{-1}SCS$ and
\begin{equation}\label{loc:three-kernel}
 K:=H'\diag(2,1)H'\diag(-3,1)H'
     =\begin{pmatrix}-8&-10\\0&-6\end{pmatrix}.
\end{equation}
Thus every upper shear with nonzero coefficient $t$ has a word with
exactly three forward copies of $C$ and invertible diagonal factors:
\begin{equation}\label{loc:nonzero-shear}
 \begin{pmatrix}1&t\\0&1\end{pmatrix}
 =D_t\diag(-1/8,-1/6)K D_t^{-1},\qquad
 D_t=\diag(4t/5,1).
\end{equation}
An update $y\gets y+t x$ is this upper shear on the ordered pair
$(y,x)$.  The order of a pair is part of the discrete schedule.

A general coefficient $\mu$ may vanish, but no test is necessary.
Prepare
\begin{equation}\label{loc:nonzero-split}
 \kappa=1+\mu\overline\mu,\qquad
 t_1=\kappa,\qquad t_2=\mu-\kappa.
\end{equation}
Both $t_1,t_2$ are nonzero, since
$1+|\mu|^2>|\mu|$.  The two corresponding shears compose to the
desired shear, including when $\mu=0$.  Equations
\eqref{loc:three-kernel}--\eqref{loc:nonzero-split} therefore give the
same six-$C$ slot pattern for every coefficient.  On disjoint pairs
these words run in parallel; every unused coordinate is an identity
singleton in every $C$ slot.

The conjugate in \eqref{loc:nonzero-split} is a prepared scalar, not
an operation on variable data.  Every coefficient $\mu$ from
Lemma~\ref{loc:fourier} is given by an arithmetic DAG over rational
numbers and the specified roots.  Evaluate a second copy of that DAG
with each root replaced by its inverse.  Rational constants are real,
so the result is precisely $\overline\mu$.  The second evaluation has
the same size, and its denominators are conjugates of nonzero
denominators.  This argument also includes the fixed root $i$, whose
conjugate is $-i$.  Shared subexpressions remain shared, so preparing
all conjugates increases polynomial preparation cost by at most a
constant factor before the constant-size replacements.

Finally, place each layer from Lemma~\ref{loc:fourier} into a
monomial slot followed by a shear slot, using identity for the other
slot type, and replace the shear slot by this fixed pattern.  Combine adjacent monomial slots if desired.  An idle round
uses identity monomials and only singleton identities in its $C$
slots.  All words now repeat a common constant-size pattern.  Padding
the shorter words with idle rounds synchronizes any finite collection
without increasing the $O(1+\log^4R)$ bound.  This proves
Proposition~\ref{loc:compiler}.

